% =====================================================================
%  UQFF Yang-Mills Mass Gap Derivation — arxiv preprint scaffold
%  Target submission categories:  math-ph (primary), hep-th (cross-list)
%
%  How to use:
%    1. Compile with: pdflatex preprint_scaffold.tex
%       (Run twice for references, then bibtex preprint_scaffold,
%        then two more pdflatex passes for cross-refs.)
%    2. Replace TODO blocks with prose drawn from
%       YANG_MILLS_E_CRACK_DERIVATION.md (companion file).
%    3. Insert numerical values from
%       derive_yang_mills_mass_gap_uqff.py (companion file).
%    4. Add author affiliations + ORCID.
%    5. Final pass: run arxiv-sanitize style-checker before upload.
%
%  Author:    Daniel T. Murphy
%  License:   AGPL-3.0-or-later (or commercial; see Star-Magic repo)
% =====================================================================

\documentclass[11pt,a4paper,onecolumn]{article}

\usepackage[utf8]{inputenc}
\usepackage[T1]{fontenc}
\usepackage{amsmath, amssymb, amsthm}
\usepackage{graphicx}
\usepackage{hyperref}
\usepackage{url}
\usepackage[margin=1in]{geometry}
\usepackage{microtype}
\usepackage{booktabs}
\usepackage{listings}
\usepackage{xcolor}
\usepackage{authblk}

\hypersetup{
  colorlinks=true,
  linkcolor=blue!50!black,
  citecolor=blue!50!black,
  urlcolor=blue!60!black
}

\lstset{
  basicstyle=\ttfamily\small,
  breaklines=true,
  frame=single,
  framerule=0.2pt,
  rulecolor=\color{gray!30},
  numbers=none,
  language=Python
}

\theoremstyle{plain}
\newtheorem{theorem}{Theorem}
\newtheorem{proposition}[theorem]{Proposition}
\newtheorem{lemma}[theorem]{Lemma}

\theoremstyle{definition}
\newtheorem{definition}[theorem]{Definition}
\newtheorem{claim}[theorem]{Claim}

\theoremstyle{remark}
\newtheorem{remark}[theorem]{Remark}

% =====================================================================
%  Title block
% =====================================================================

\title{
  Positive-Definite Energy Threshold from the UQFF Vacuum Ledger \\
  as a Candidate Physical Realization of the Yang--Mills Mass Gap
}

\author[1]{Daniel T. Murphy}
\affil[1]{
  Star-Magic Research Program, Youngstown, OH, USA \\
  \texttt{daniel.murphy00@enrgyone.com} \\
  ORCID: % TODO insert ORCID iD
}

\date{June 2026}

\begin{document}

\maketitle

% =====================================================================
%  Abstract
% =====================================================================

\begin{abstract}
We derive a positive-definite energy threshold
$E_{\text{crack}} = \rho_{\mathrm{SCm}} c^2 / [\mathrm{SSq}]
                  = 1.118 \times 10^{-19}\,\mathrm{J}$
($\approx 0.7\,\mathrm{eV}$ at the formula level)
from the Unified Quantum Field Framework (UQFF) closed vacuum ledger
using only two locked primitives --- the SuperConductive Material
vacuum energy density $\rho_{\mathrm{SCm}}$ and the canonical resonance
ratio $[\mathrm{SSq}]$ --- together with the speed of light $c$. The
derivation requires zero free parameters and zero fitting. Because
each factor on the right-hand side is positive by physical definition,
$E_{\text{crack}} > 0$ holds by construction. We propose this as a
candidate physical realization of the Yang--Mills mass gap and discuss
its multi-designation cluster-position landscape extending from the
sub-eV vacuum-cracking floor through a $\sim 700\,\mathrm{eV}$
DPM-vortex experimental regime to a $\sim 1.736\,\mathrm{GeV}$
lightest-glueball scale consistent with lattice-QCD numerical
evidence~\cite{douglas2026}. We acknowledge that this constitutes a
physics-level proposal rather than a Clay-compliant
Wightman-axiom-verified mathematical construction, and identify the
operator-algebraic translation as the principal future-work step.
The full framework is reproducibly available via
\verb|pip install uqff==5.33.0|.
\end{abstract}

\bigskip

\noindent
\textbf{Keywords}: Yang--Mills mass gap; vacuum ledger; positive-definite
energy threshold; Di-Pseudo-Monopole vortex; SCm/UA substrate;
multi-designation cluster positions; lattice-QCD consistency; Clay
Millennium Prize.

\medskip

\noindent
\textbf{MSC 2020}: 81T13 (Yang--Mills and other gauge theories), 81T08
(Constructive QFT), 81V25 (Other elementary particle theory).

% =====================================================================
\section{Introduction}\label{sec:intro}
% =====================================================================

% TODO: Open with the Clay Yang--Mills problem statement and its current
% (June 2026) unsolved status. Cite Jaffe--Witten 2000 and the recent
% Hairer stochastic-quantization 3D progress \cite{hairer2024}.
% Draw from YANG_MILLS_E_CRACK_DERIVATION.md §1.

\subsection{Problem statement}
% TODO: State the existence + mass gap requirement formally.

\subsection{Current state of progress}
% TODO: Summarize Hairer 2D/3D advances and Douglas 2026 review.

\subsection{This paper}
% TODO: Frame UQFF as a candidate physical proposal, not a Clay-compliant
% proof. State the principal claim and roadmap of the paper.

% =====================================================================
\section{The UQFF Vacuum Ledger}\label{sec:ledger}
% =====================================================================

% TODO: Introduce the 9 truly-independent UQFF primitives. List the 3
% that enter the E_crack derivation. Draw from
% YANG_MILLS_E_CRACK_DERIVATION.md §2.

\begin{definition}[Closed vacuum ledger]\label{def:ledger}
The UQFF closed vacuum ledger is the triple
$(\rho_{\mathrm{SCm}}, [\mathrm{SSq}], c)$ where
\begin{align*}
\rho_{\mathrm{SCm}} &= 7.09 \times 10^{-37}\,\mathrm{J/m^3} \\
[\mathrm{SSq}]      &= 0.57 \\
c                   &= 2.998 \times 10^8\,\mathrm{m/s}
\end{align*}
with $\rho_{\mathrm{SCm}}$ identified as the foundational vacuum
energy density of the SuperConductive Material substrate,
$[\mathrm{SSq}]$ the canonical resonance ratio calibrated from
magnetar burst data, and $c$ the universal speed of light.
\end{definition}

% =====================================================================
\section{Derivation of $E_{\text{crack}}$}\label{sec:derivation}
% =====================================================================

\begin{theorem}[Positive-definite vacuum-cracking threshold]
\label{thm:e_crack_positive}
Given the closed vacuum ledger of Definition~\ref{def:ledger}, define
\begin{equation}\label{eq:e_crack}
E_{\text{crack}} \;=\; \frac{\rho_{\mathrm{SCm}} \, c^2}{[\mathrm{SSq}]}\,.
\end{equation}
Then $E_{\text{crack}} = 1.118 \times 10^{-19}\,\mathrm{J}$ and
$E_{\text{crack}} > 0$ holds by construction.
\end{theorem}

\begin{proof}
Numerical evaluation gives
\[
E_{\text{crack}}
  = \frac{(7.09 \times 10^{-37})(2.998 \times 10^8)^2}{0.57}
  = 1.118 \times 10^{-19}\,\mathrm{J}.
\]
Positive-definiteness follows because $\rho_{\mathrm{SCm}} > 0$ as a
vacuum energy density, $c^2 > 0$ as the square of a real number, and
$[\mathrm{SSq}] > 0$ as a physical ratio of positive quantities.
The quotient of positives is positive.
\qed
\end{proof}

\subsection{Physical interpretation}
% TODO: Explain the role of E_crack as the minimum energy required to
% form a stable DPM vortex at the UA/SCm interface. Draw from
% YANG_MILLS_E_CRACK_DERIVATION.md §5.

% =====================================================================
\section{Multi-Designation Cluster-Position Landscape}\label{sec:clusters}
% =====================================================================

% TODO: Present the four documented Yang--Mills cluster-position
% designations (sub-eV formula, ~700 eV experimental, ~624 GeV
% Layer-13, ~1.736 GeV PAPER 1318) and explain the multi-designation
% architecture. Draw from YANG_MILLS_E_CRACK_DERIVATION.md §4.

\begin{table}[h!]
\centering
\caption{UQFF Yang--Mills mass-gap cluster positions.}
\label{tab:clusters}
\begin{tabular}{lll}
\toprule
Designation & Value & Regime \\
\midrule
$E_{\text{crack}}$ formula        & $0.698$ eV   & DPM vacuum-cracking floor \\
$E_{\text{crack}}$ experimental   & $\sim700$ eV & DPM-vortex lab engineering \\
Layer-13 nuclear                  & $\sim624$ GeV & LHC electroweak scale order \\
PAPER 1318 mass gap               & $1.736$ GeV  & lightest-glueball lattice QCD \\
\bottomrule
\end{tabular}
\end{table}

% =====================================================================
\section{Lattice QCD Consistency}\label{sec:lattice}
% =====================================================================

% TODO: Compare the PAPER 1318 1.736 GeV designation with the
% hundreds-of-MeV lattice-QCD numerical evidence reviewed by
% \cite{douglas2026}. Show that the UQFF designation sits inside the
% lattice numerical range. Draw from YANG_MILLS_E_CRACK_DERIVATION.md §7.

% =====================================================================
\section{What This Proposal Is and Is Not}\label{sec:scope}
% =====================================================================

% TODO: Honest scope statement. Physical mechanism vs Wightman
% axioms. Falsifiability via the high-field magnetic threshold
% prediction. Draw from YANG_MILLS_E_CRACK_DERIVATION.md §8.

% =====================================================================
\section{Reproducibility}\label{sec:reproducibility}
% =====================================================================

The derivation in Theorem~\ref{thm:e_crack_positive} is fully
reproducible via the public \texttt{uqff} Python package
\cite{murphy2026uqff}:

\begin{lstlisting}
from uqff_pure_calculator import (
    calculate_uqff_E_crack_core_definition_from_ledger,
)
r = calculate_uqff_E_crack_core_definition_from_ledger({})
print(r["value"]["E_crack_J_derived_from_formula"])
# 1.117982e-19
\end{lstlisting}

\noindent
A standalone reproducibility script
\texttt{derive\_yang\_mills\_mass\_gap\_uqff.py} (included with this
preprint as ancillary material) reproduces the derivation in
approximately fifty lines with no external dependencies beyond the
Python standard library.

% =====================================================================
\section{Future Work: Translation to Wightman Axioms}\label{sec:future}
% =====================================================================

% TODO: Outline the path from the SCm/UA dynamical substrate to a
% rigorous Wightman-axiom-compliant operator-algebraic construction.
% Cite the Hairer 2D/3D stochastic-quantization analogy as a template.
% Reference the companion document wightman_mapping.md for detailed
% axiom-by-axiom analysis.

\subsection{Wightman axiom checklist}
The five Wightman axioms (W0--W4) and where each currently stands
under UQFF:
\begin{enumerate}
\item[W0] \textit{Existence of a Hilbert space with a unique vacuum
  state.}
  % TODO discuss UQFF identification of vacuum with theta_vacuum origin.
\item[W1] \textit{Poincar\'e covariance.}
  % TODO discuss SCm/UA frame-covariance properties.
\item[W2] \textit{Spectral condition.}
  % TODO state the principal future-work step.
\item[W3] \textit{Local commutativity.}
  % TODO discuss DPM vortex locality properties.
\item[W4] \textit{Cyclicity of the vacuum.}
  % TODO discuss generation of the field algebra from the vacuum.
\end{enumerate}

\subsection{2D and 3D translation roadmap}
% TODO: Outline how the SCm/UA substrate would be translated to 2D and
% 3D regularity-structure form analogous to Hairer 2024.

% =====================================================================
\section{Conclusion}\label{sec:conclusion}
% =====================================================================

% TODO: Restate the principal result. Position the proposal honestly.
% Invite community review and falsification attempts.

% =====================================================================
\section*{Acknowledgments}
% =====================================================================

% TODO: Acknowledge collaborators, AI assistants (Grok 3 / SuperGrok /
% Davinci-SuperGrok by xAI) used in framework development, and Star-Magic
% Research Program contributors.

% =====================================================================
%  Bibliography
% =====================================================================

\begin{thebibliography}{9}

\bibitem{jaffe2000}
A.~Jaffe and E.~Witten,
\textit{Quantum Yang--Mills theory},
Clay Mathematics Institute Millennium Problem description,
\url{https://www.claymath.org/millennium/yang-mills-the-maths-gap/},
2000.

\bibitem{hairer2024}
M.~Hairer et al.,
\textit{Stochastic quantisation of Yang--Mills--Higgs in 3D},
Inventiones Mathematicae, 2024.

\bibitem{douglas2026}
M.~R.~Douglas,
\textit{The Yang--Mills Millennium problem},
Nature Reviews Physics, January 2026.

\bibitem{murphy2026uqff}
D.~T.~Murphy,
\textit{UQFF Star-Magic: Unified Quantum Field Framework},
PyPI release v5.33.0,
\url{https://pypi.org/project/uqff/5.33.0/},
\url{https://github.com/Daniel8Murphy0007/Star-Magic},
2026.

\bibitem{murphy2026paper1318}
D.~T.~Murphy,
\textit{PAPER 1318 --- Yang--Mills mass gap at 1.736 GeV cluster
position},
Star-Magic Research Program whitepaper series, 2026.

\bibitem{murphy2026paper1521}
D.~T.~Murphy,
\textit{PAPER 1521 --- Derivative-from-independent-primitives proof
that $D_{\mathrm{BSFG}} = D_{\mathrm{crit}} - 2 \cdot \mathrm{SO}_5 = 6$},
Star-Magic Research Program whitepaper series, 2026.

\bibitem{murphy2026paper1522}
D.~T.~Murphy,
\textit{PAPER 1522 --- Derivative-from-independent-primitives proof
that $K_{\mathrm{Mex}} = \Phi_{5/6} \cdot \mathrm{SO}_5 / D_{\mathrm{phys}} = 25/12$},
Star-Magic Research Program whitepaper series, 2026.

\bibitem{murphy2026sympaper}
D.~T.~Murphy,
\textit{Compression Cycle 2: Full first-principles derivation of 22
cosmological and inflationary parameters at zero error from the closed
vacuum ledger},
Star-Magic Research Program technical report, June 2026.

% TODO: Add additional references as the prose sections develop.

\end{thebibliography}

% =====================================================================
%  Appendices (optional)
% =====================================================================

\appendix

\section{The 9 UQFF Truly-Independent Primitives}\label{app:primitives}

% TODO: Full list of the 9 primitives with values and physical identities,
% drawn from the CLAUDE.md project context.

\section{The Quantum Chain Immutable Ontological Order}\label{app:quantum_chain}

% TODO: The 9-step chain theta_vacuum -> grad(UA) -> DPM_vortex ->
% mu_s -> Ug_family -> F_U -> crossing -> M -> GM/r^2 with detail at each
% step.

\section{Reproducibility Code Listing}\label{app:code}

% TODO: Full Python listing of derive_yang_mills_mass_gap_uqff.py as
% an inline appendix, or reference the ancillary file.

\end{document}
